% GENERATED FILE. Edit canonical lessons, then run npm run generate:books.
% canonical-source-hash: fnv1a64:be3e705f68b764f8
% canonical-lessons: SA-C117-angusthah, SA-C117-vaksah, SA-C117-katih, SA-C117-bhruh, SA-C117-tvak

\chapter{Thumb, Chest, Waist, Eyebrow, Skin}
\label{ch:sa-thumb-chest-waist-eyebrow-skin}
\hlchaptermodality{117}

\begin{chapteropening}
\textbf{By the end of this chapter:} \emph{I can say a thumb, the chest, the waist, an eyebrow and the skin.}

Complete the last lesson of chapter 117: i can say a thumb, the chest, the waist, an eyebrow and the skin.
\end{chapteropening}

\section[aṅguṣṭhaḥ]{\sk{अंगुष्ठः} (aṅguṣṭhaḥ) — a thumb}

\label{lesson:SA-C117-angusthah}



\begin{quote}
Before the new one: say the Sanskrit for an elbow, then the Sanskrit for a wrist.
\end{quote}

\subsection*{You'll want to know: \sk{अंगुष्ठः}}
\textbf{\sk{अंगुष्ठः}} — \emph{aṅguṣṭhaḥ} — \textquotedblleft{}a thumb\textquotedblright{}.

\begin{grammarlens}[title={What you've built}]
One more word for this chapter.
\end{grammarlens}

\subsection*{Your turn}
\begin{itemize}
\raggedright
  \item \emph{Say it:} \emph{aṅguṣṭhaḥ}
  \item \emph{Say it:} \emph{aṅguṣṭhaḥ}, once more
  \item \emph{Say it:} say \emph{aṅguṣṭhaḥ}
  \item \emph{Recall:} say the Sanskrit for an elbow, then the Sanskrit for a wrist, then say \emph{aṅguṣṭhaḥ} again
\end{itemize}

\begin{tcolorbox}[breakable,colback=teal!4,colframe=teal!35!black,title={Before you move on}]
What does \sk{अंगुष्ठः} mean? (\textquotedblleft{}A thumb\textquotedblright{}.) Say it once more.
\end{tcolorbox}
\section[vakṣaḥ]{\sk{वक्षः} (vakṣaḥ) — the chest}

\label{lesson:SA-C117-vaksah}



\begin{quote}
Before the new one: say the Sanskrit for a wrist, then the Sanskrit for a thumb.
\end{quote}

\subsection*{You'll want to know: \sk{वक्षः}}
\textbf{\sk{वक्षः}} — \emph{vakṣaḥ} — \textquotedblleft{}the chest\textquotedblright{}.

\begin{grammarlens}[title={What you've built}]
One more word for this chapter.
\end{grammarlens}

\subsection*{Your turn}
\begin{itemize}
\raggedright
  \item \emph{Say it:} \emph{vakṣaḥ}
  \item \emph{Say it:} \emph{vakṣaḥ}, once more
  \item \emph{Say it:} say \emph{vakṣaḥ}
  \item \emph{Recall:} say the Sanskrit for a wrist, then the Sanskrit for a thumb, then say \emph{vakṣaḥ} again
\end{itemize}

\begin{tcolorbox}[breakable,colback=teal!4,colframe=teal!35!black,title={Before you move on}]
What does \sk{वक्षः} mean? (\textquotedblleft{}The chest\textquotedblright{}.) Say it once more.
\end{tcolorbox}
\section[kaṭiḥ]{\sk{कटिः} (kaṭiḥ) — the waist}

\label{lesson:SA-C117-katih}



\begin{quote}
Before the new one: say the Sanskrit for a thumb, then the Sanskrit for the chest.
\end{quote}

\subsection*{You'll want to know: \sk{कटिः}}
\textbf{\sk{कटिः}} — \emph{kaṭiḥ} — \textquotedblleft{}the waist\textquotedblright{}.

\begin{grammarlens}[title={What you've built}]
One more word for this chapter.
\end{grammarlens}

\subsection*{Your turn}
\begin{itemize}
\raggedright
  \item \emph{Say it:} \emph{kaṭiḥ}
  \item \emph{Say it:} \emph{kaṭiḥ}, once more
  \item \emph{Say it:} say \emph{kaṭiḥ}
  \item \emph{Recall:} say the Sanskrit for a thumb, then the Sanskrit for the chest, then say \emph{kaṭiḥ} again
\end{itemize}

\begin{tcolorbox}[breakable,colback=teal!4,colframe=teal!35!black,title={Before you move on}]
What does \sk{कटिः} mean? (\textquotedblleft{}The waist\textquotedblright{}.) Say it once more.
\end{tcolorbox}
\section[bhrūḥ]{\sk{भ्रूः} (bhrūḥ) — an eyebrow}

\label{lesson:SA-C117-bhruh}



\begin{quote}
Before the new one: say the Sanskrit for the chest, then the Sanskrit for the waist.
\end{quote}

\subsection*{You'll want to know: \sk{भ्रूः}}
\textbf{\sk{भ्रूः}} — \emph{bhrūḥ} — \textquotedblleft{}an eyebrow\textquotedblright{}.

\begin{grammarlens}[title={What you've built}]
One more word for this chapter.
\end{grammarlens}

\subsection*{Your turn}
\begin{itemize}
\raggedright
  \item \emph{Say it:} \emph{bhrūḥ}
  \item \emph{Say it:} \emph{bhrūḥ}, once more
  \item \emph{Say it:} say \emph{bhrūḥ}
  \item \emph{Recall:} say the Sanskrit for the chest, then the Sanskrit for the waist, then say \emph{bhrūḥ} again
\end{itemize}

\begin{tcolorbox}[breakable,colback=teal!4,colframe=teal!35!black,title={Before you move on}]
What does \sk{भ्रूः} mean? (\textquotedblleft{}An eyebrow\textquotedblright{}.) Say it once more.
\end{tcolorbox}
\section[tvak]{\sk{त्वक्} (tvak) — the skin}

\label{lesson:SA-C117-tvak}



\begin{quote}
Before the new one: say the Sanskrit for the waist, then the Sanskrit for an eyebrow.
\end{quote}

\subsection*{You'll want to know: \sk{त्वक्}}
\textbf{\sk{त्वक्}} — \emph{tvak} — \textquotedblleft{}the skin\textquotedblright{}.

\begin{grammarlens}[title={What you've built}]
One more word for this chapter.
\end{grammarlens}

\subsection*{Your turn}
\begin{itemize}
\raggedright
  \item \emph{Say it:} \emph{tvak}
  \item \emph{Say it:} \emph{tvak}, once more
  \item \emph{Say it:} say \emph{tvak}
  \item \emph{Recall:} say the Sanskrit for the waist, then the Sanskrit for an eyebrow, then say \emph{tvak} again
\end{itemize}

\begin{tcolorbox}[breakable,colback=teal!4,colframe=teal!35!black,title={Before you move on}]
What does \sk{त्वक्} mean? (\textquotedblleft{}The skin\textquotedblright{}.) Say it once more.
\end{tcolorbox}
